\section{有效性威胁}
\label{sec:threats}

我们的数字能分辨的差异大约是 30 道题中的两到三道，所以我们能下的结论是：我们上线或否决的任何单个手段，都没有让留出
得分变化超过这个幅度。

\begin{itemize}[leftmargin=1.4em]
  \item \textbf{噪声。}在未读的 30 道题上，两次对两次的读取，零假设标准差约为 1.4；我们在那里测量的每个单变量手段，都落在
    基线的 2 以内；只有第~\ref{sec:structure}~节的配置变化（人设与推理强度一起变）超出了这个范围。更小的真实效应，无论正负，
    都显示不出来。
  \item \textbf{在同一留出集上做了很多检验。}未读的 30 道题从未打开，但有几十个手段在它们上面读取过结果，每次读取都会
    泄露一点信息 \citep{dwork2015reusable}。一个靠运气通过的手段可能被误采用；我们采用的手段很少，主要的那一个经过了
    多次运行的确认。
  \item \textbf{抽样受额度驱动。}运行来自一个有每周额度的订阅账号。额度不足时实验暂停（10 月 26 日、11 月 8 日、
    11 月 20 日），筛查先在更便宜的模型上跑，一些实验组在已经不可能达到门槛时提前停止。
  \item \textbf{账号登录，而非 API。}几乎所有运行都使用 ChatGPT（Codex）账号登录。成本是根据服务端价格表按公开标价
    重新计算的；API 路径只通过 OpenRouter 做过一次四次尝试的冒烟测试。API 上的速率限制、延迟和缓存可能不同。
  \item \textbf{只有一台机器。}30 次尝试共用一台台式电脑，所以比在每题独占的云机器上慢；验证器在负载下有时会超时
    （由空闲重评处理）。这台电脑没有 GPU，所以 66 道题中有 3 道从未运行。
  \item \textbf{自跑的结果。}Terminal-Bench~4.1 的运行和评分都由我们自己完成，而不是由 AA 或 Terminal-Bench 维护者完成。
    适配器已经公开，以便复现（附录~\ref{app:repro}）。
  \item \textbf{两边都是我们做的。}产品、V1 任务集及其评分器都是我们的。盲做抓出了两个评分器缺陷，之后的审计也撤下或
    修复了不公平的题，但可能还有遗漏。我们的奖励投机扫描是启发式的，它标出的两次留出尝试尚未判断。
  \item \textbf{产品在变。}产品在实验期间一直在变，这正是每次比较都在同一批次内进行的原因。GPT-6.1 上的结论未必适用于
    GPT-6.2，Terminal-Bench~4.0 上的结论也未必适用于 4.1（有一道题不同）。
\end{itemize}
